\clearpage
\chapter*{2023 Septiembre}
\addcontentsline{toc}{chapter}{2023 Septiembre}

\begin{enunciado}
    \hypertarget{2023 S 1}{}
    \textbf{1. (3,75 puntos)} Sean los espacios \((\mathbb{R},T_u)\) e \((I,T)\), con \(I=[0,1]\subset \mathbb{R}\) y \(T\) la topología de \(I\) como subespacio de \((\mathbb{R},T_u)\). Se considera la aplicación \(f:(\mathbb{R},T_u)\to (I,T)\) definida por
    \[
        f(x) = \begin{cases}
            0 & \text{si } x \leq 0 \\
            x & \text{si } 0 \leq x \leq 1 \\
            1 & \text{si } x \geq 1
        \end{cases}
    \]
    \begin{enumerate}[label=\alph*)]
        \item Estudiar si \(f\) es una aplicación continua.
        \item Estudiar si \(f\) es abierta.
        \item Estudiar si \(f\) es cerrada.
    \end{enumerate}
\end{enunciado}

\begin{enunciado}
    \hypertarget{2023 S 2}{}
    \textbf{2. (3,75 puntos)} Dado el conjunto \(\mathbb{R}\) de los números reales, se considera la topología \(T = \{\emptyset\} \cup \{G \subset \mathbb{R} \mid 0 \in G\}\).
    \begin{enumerate}[label=\alph*)]
        \item Estudiar si \((\mathbb{R},T)\) verifica los axiomas I y II de numerabilidad.
        \item Determinar las sucesiones en \((\mathbb{R},T)\) que convergen a 1. ¿Es único dicho límite?
        \item Estudiar si \((\mathbb{R},T)\) es un espacio topológico conexo.
    \end{enumerate}
\end{enunciado}

\begin{enunciado}
    \hypertarget{2023 S 3}{}
    \textbf{3. (2,5 puntos)}
    \begin{enumerate}[label=\alph*)]
        \item Sea \(M\) un subconjunto de un espacio topológico \((X,T)\). Dar la definición de punto de acumulación de \(M\).
        \item Sea \((X,T)\) un espacio topológico compacto, y sea \(M\) un subconjunto infinito de \(X\). Estudiar si existe necesariamente un punto de acumulación de \(M\) en \(X\).
    \end{enumerate}
\end{enunciado}